\documentclass[11pt]{article}
% Preámbulo compartido por los artefactos Scrum del proyecto
% Proyecto Agile — Spanish Braille Application (EPN)
\usepackage[utf8]{inputenc}
\usepackage[T1]{fontenc}
\usepackage[spanish]{babel}
\usepackage[a4paper,margin=2.3cm]{geometry}
\usepackage{amsmath}
\usepackage{amssymb}
\usepackage{xcolor}
\usepackage{booktabs}
\usepackage{longtable}
\usepackage{array}
\usepackage{colortbl}
\usepackage{enumitem}
\usepackage{fancyhdr}
\usepackage{titlesec}
\usepackage{tcolorbox}
\usepackage{pgfplots}
\pgfplotsset{compat=1.18}
\usepackage{hyperref}

\definecolor{epnblue}{HTML}{2E4C9C}
\definecolor{epnbluelight}{HTML}{4B6FCB}
\definecolor{epnrow}{HTML}{EDF1FB}
\definecolor{epngray}{HTML}{555555}

\hypersetup{
  colorlinks=true,
  linkcolor=epnblue,
  urlcolor=epnblue,
  citecolor=epnblue,
  pdfcreator={XeLaTeX},
}

\titleformat{\section}
  {\normalfont\Large\bfseries\color{epnblue}}{\thesection}{0.6em}{}
\titleformat{\subsection}
  {\normalfont\large\bfseries\color{epnbluelight}}{\thesubsection}{0.6em}{}

\pagestyle{fancy}
\fancyhf{}
\renewcommand{\headrulewidth}{0.4pt}
\renewcommand{\footrulewidth}{0.2pt}
\fancyfoot[C]{\thepage}

\newcommand{\artefactoTitulo}[2]{%
  \begin{titlepage}
    \centering
    \vspace*{1.2cm}
    {\large Escuela Politécnica Nacional\par}
    {\normalsize Facultad de Ingeniería de Sistemas\par}
    {\normalsize Construcción y Evolución de Software --- Primer Bimestre\par}
    \vspace{1.4cm}
    \begin{tcolorbox}[colback=epnblue,colframe=epnblue,coltext=white,arc=2mm,boxrule=0pt,width=0.92\textwidth,halign=center]
      {\Huge\bfseries #1\par}
    \end{tcolorbox}
    \vspace{0.5cm}
    {\Large #2\par}
    \vspace{1.6cm}
    {\large\bfseries Proyecto Agile --- Spanish Braille Application\par}
    \vspace{0.3cm}
    {\normalsize Traducción bidireccional español \(\leftrightarrow\) Braille Unicode\par}
    \vfill
    \begin{tabular}{ll}
      \textbf{Product Owner:}  & José Castro \\
      \textbf{Scrum Master:}   & Victor Aveiga \\
      \textbf{Equipo de Desarrollo:} & Javier Angulo, Elian Caizapanta, Erick Costa, Emily Aumala \\
    \end{tabular}
    \vspace{0.8cm}
    \par\noindent\rule{0.5\textwidth}{0.4pt}\par
    \vspace{0.3cm}
    {\normalsize Documento preparado por: Erick Costa (Equipo de Desarrollo)\par}
    {\normalsize \today\par}
  \end{titlepage}
  \fancyhead[L]{Proyecto Agile --- Spanish Braille App}
  \fancyhead[R]{#1}
}


\begin{document}
\artefactoTitulo{Sprint Backlog}{Desglose de tareas técnicas por sprint y responsable}

\section{Sprint 1 --- Transcripción Base}
\textbf{Objetivo del Sprint:} implementar la conversión completa de texto español a Braille, incluyendo
abecedario, vocales acentuadas, números y mayúsculas.
\quad\textbf{Fechas reales:} 22 jun -- 3 jul 2026 (2 semanas)
\quad\textbf{Capacidad:} 12 pts \quad \textbf{HUs:} 01H1, 02H2, 03H3, 04H4

\renewcommand{\arraystretch}{1.3}
\begin{longtable}{|p{3cm}|p{1cm}|p{1.1cm}|p{7.6cm}|}
\hline
\rowcolor{epnblue}
\textcolor{white}{\textbf{Responsable}} & \textcolor{white}{\textbf{HU}} & \textcolor{white}{\textbf{Horas}} & \textcolor{white}{\textbf{Tarea técnica}} \\
\hline
\endfirsthead
\hline
\rowcolor{epnblue}
\textcolor{white}{\textbf{Responsable}} & \textcolor{white}{\textbf{HU}} & \textcolor{white}{\textbf{Horas}} & \textcolor{white}{\textbf{Tarea técnica}} \\
\hline
\endhead
\rowcolor{epnrow}
Erick Costa & 01H1 & 5h & Mapa de caracteres del abecedario español (series 1.\textsuperscript{a}/2.\textsuperscript{a}/3.\textsuperscript{a}) \\
\hline
Emily Aumala & 01H1 & 4h & Función de conversión letra $\rightarrow$ patrón Braille \\
\hline
\rowcolor{epnrow}
Javier Angulo & --- & 3h & Interfaz de entrada de texto (\texttt{index.html}) \\
\hline
Javier Angulo & 02H2 & 4h & Extensión del mapa de caracteres: vocales acentuadas y ñ \\
\hline
\rowcolor{epnrow}
Elian Caizapanta & 02H2 & 3h & Pruebas de integración de vocales acentuadas y ñ \\
\hline
Erick Costa & 03H3 & 4h & Detección de dígitos y signo numérico Braille (puntos 3-4-5-6) \\
\hline
\rowcolor{epnrow}
Emily Aumala & 03H3 & 3h & Validación de secuencias numéricas multi-cifra \\
\hline
Javier Angulo & 04H4 & 3h & Detección de mayúsculas e indicador Braille (puntos 4-6) \\
\hline
\rowcolor{epnrow}
Elian Caizapanta & 04H4 & 3h & Pruebas de integración de mayúsculas y enrutamiento del backend \\
\hline
\end{longtable}
\textbf{Total Sprint 1:} 32 h (Erick Costa 9h, Emily Aumala 7h, Javier Angulo 10h, Elian Caizapanta 6h)

\section{Sprint 2 --- Funcionalidades Avanzadas}
\textbf{Objetivo del Sprint:} agregar puntuación, generar señalética imprimible y habilitar la
transcripción inversa Braille $\rightarrow$ español.
\quad\textbf{Fechas reales:} 6 -- 10 jul 2026 (1 semana)
\quad\textbf{Capacidad:} 15 pts \quad \textbf{HUs:} 05H5, 06H6, 07H7

\begin{longtable}{|p{3cm}|p{1cm}|p{1.1cm}|p{7.6cm}|}
\hline
\rowcolor{epnblue}
\textcolor{white}{\textbf{Responsable}} & \textcolor{white}{\textbf{HU}} & \textcolor{white}{\textbf{Horas}} & \textcolor{white}{\textbf{Tarea técnica}} \\
\hline
\endfirsthead
\hline
\rowcolor{epnblue}
\textcolor{white}{\textbf{Responsable}} & \textcolor{white}{\textbf{HU}} & \textcolor{white}{\textbf{Horas}} & \textcolor{white}{\textbf{Tarea técnica}} \\
\hline
\endhead
\rowcolor{epnrow}
Erick Costa & 05H5 & 3h & Mapa de signos de puntuación básicos (\texttt{, . ; : ¿ ? ¡ !}) \\
\hline
Emily Aumala & 05H5 & 2h & Pruebas de oraciones completas con signos de puntuación \\
\hline
\rowcolor{epnrow}
Javier Angulo & 06H6 & 4h & Plantilla de señalética Braille imprimible (Braille espejo) \\
\hline
Elian Caizapanta & 06H6 & 5h & Exportación de señalética (renderizado a formato descargable) \\
\hline
\rowcolor{epnrow}
Erick Costa & 07H7 & 5h & Función inversa Braille $\rightarrow$ español (diccionario inverso e indicadores) \\
\hline
Emily Aumala & 07H7 & 4h & Interfaz de entrada de patrones Braille para la conversión inversa \\
\hline
\end{longtable}
\textbf{Total Sprint 2:} 23 h (Erick Costa 5h, Emily Aumala 6h, Javier Angulo 4h, Elian Caizapanta 5h)

\vspace{0.2cm}
\noindent\textit{Nota sobre 06H6:} el equipo reporta la tarea de exportación de señalética (Elian
Caizapanta) como completada. Al momento de preparar este documento, la rama clonada de este repositorio
(\texttt{Proyecto-Agile}) todavía no refleja un endpoint de exportación PNG/PDF dedicado --- solo la vista
de Braille espejo con impresión de navegador. Se recomienda confirmar que el código esté sincronizado
con la rama \texttt{ITERACION-2} del repositorio de código fuente antes de la entrega final, para que el
Incremento demostrado coincida con lo documentado aquí.

\section{Horas totales por integrante (ambos sprints)}
\begin{center}
\renewcommand{\arraystretch}{1.2}
\begin{tabular}{|l|c|}
\hline
\rowcolor{epnblue}
\textcolor{white}{\textbf{Integrante}} & \textcolor{white}{\textbf{Horas técnicas}} \\
\hline
Javier Angulo & 14h \\
\hline
\rowcolor{epnrow}
Elian Caizapanta & 12h \\
\hline
Erick Costa & 14h \\
\hline
\rowcolor{epnrow}
Emily Aumala & 11h \\
\hline
\end{tabular}
\end{center}

\section{Convenciones de trabajo}
\begin{itemize}[leftmargin=1.4cm]
  \item El \textbf{Sprint Backlog} es un artefacto creado y de propiedad del propio \textbf{Equipo de
        Desarrollo} durante el Sprint Planning; nadie externo (ni el Scrum Master ni el Product Owner)
        asigna las tareas --- el equipo se autoorganiza y cada integrante toma tareas según disponibilidad
        y habilidad.
  \item Cada sprint es un \textbf{time-box} de duración fija acordada de antemano (Sprint 1: 2 semanas;
        Sprint 2: 1 semana), cerrado en la fecha planificada con un objetivo propio, sin extender el plazo
        aunque quedaran ítems abiertos (ver nota sobre 06H6).
  \item El \textbf{Scrum Master (Victor Aveiga)} facilita el Daily Scrum (15 registros diarios entre ambos
        sprints) y remueve impedimentos; el \textbf{Product Owner (José Castro)} prioriza el Product
        Backlog y acepta el incremento en cada Sprint Review.
\end{itemize}

\end{document}
